\chapter{Activation Functions and Architectures}

\section{Nonlinear Activations}

Nonlinear activation functions are a fundamental component of neural networks. They enable networks to represent complex, non-linear relationships that cannot be captured by linear transformations alone.

\medskip

\noindent
\textbf{Role of activation functions.}  
Each layer in a neural network applies a transformation of the form
\[
z \mapsto \sigma(W z + b),
\]
where $\sigma$ is applied elementwise.

\medskip

\noindent
\textbf{Why nonlinearity is necessary.}  
If $\sigma$ is linear, then a composition of layers reduces to a single linear transformation:
\[
W_L W_{L-1} \cdots W_1 x,
\]
eliminating the benefit of depth.

\medskip

\noindent
\textbf{Common activation functions.}

\begin{itemize}
    \item \textbf{Sigmoid:}
    \[
    \sigma(z) = \frac{1}{1 + e^{-z}}
    \]

    \item \textbf{Tanh:}
    \[
    \sigma(z) = \tanh(z)
    \]

    \item \textbf{ReLU (Rectified Linear Unit):}
    \[
    \sigma(z) = \max(0, z)
    \]
\end{itemize}

\medskip

\noindent
\textbf{Properties.}
\begin{itemize}
    \item Sigmoid and tanh are smooth but can saturate
    \item ReLU is simple and avoids saturation in the positive region
\end{itemize}

\medskip

\noindent
\textbf{Variants.}
\begin{itemize}
    \item Leaky ReLU
    \item ELU
    \item GELU
\end{itemize}

\medskip

\noindent
\textbf{Gradient behavior.}
\begin{itemize}
    \item Saturating activations $\rightarrow$ vanishing gradients
    \item Piecewise linear activations $\rightarrow$ stable gradients
\end{itemize}

\medskip

\noindent
\textbf{Insight.}  
The choice of activation function strongly influences both expressivity and optimization dynamics.

\section{Architectural Design Patterns}

Beyond individual layers, the overall structure of a neural network plays a crucial role in its performance.

\medskip

\noindent
\textbf{Fully connected networks.}  
Each layer connects to all units in the previous layer.

\medskip

\noindent
\textbf{Convolutional architectures.}
\begin{itemize}
    \item Exploit spatial locality
    \item Share parameters across locations
\end{itemize}

\medskip

\noindent
\textbf{Residual connections.}
\begin{itemize}
    \item Add identity mappings between layers:
    \[
    z^{(\ell+1)} = z^{(\ell)} + F(z^{(\ell)})
    \]
    \item Improve gradient flow and stability
\end{itemize}

\medskip

\noindent
\textbf{Normalization layers.}
\begin{itemize}
    \item Batch normalization
    \item Layer normalization
\end{itemize}

\medskip

\noindent
\textbf{Modular design.}  
Modern architectures are often composed of repeated blocks (e.g., convolutional blocks, transformer blocks).

\medskip

\noindent
\textbf{Insight.}  
Architectural design encodes inductive biases that guide learning.

\section{Depth versus Width}

The structure of a neural network is characterized by its depth and width.

\medskip

\noindent
\textbf{Width.}  
Number of units in each layer.

\medskip

\noindent
\textbf{Depth.}  
Number of layers.

\medskip

\noindent
\textbf{Shallow wide networks.}
\begin{itemize}
    \item Can approximate many functions
    \item May require large numbers of parameters
\end{itemize}

\medskip

\noindent
\textbf{Deep networks.}
\begin{itemize}
    \item Capture hierarchical structure
    \item Often more parameter-efficient
\end{itemize}

\medskip

\noindent
\textbf{Expressivity comparison.}  
Certain functions can be represented efficiently by deep networks but require exponentially large shallow networks.

\medskip

\noindent
\textbf{Optimization considerations.}
\begin{itemize}
    \item Deeper networks are harder to train
    \item Techniques such as residual connections mitigate this
\end{itemize}

\medskip

\noindent
\textbf{Insight.}  
Depth provides a natural way to encode compositional structure.

\section{Expressivity Considerations}

The expressive power of a neural network depends on its architecture and activation functions.

\medskip

\noindent
\textbf{Expressivity.}  
The ability of a model to approximate a wide range of functions.

\medskip

\noindent
\textbf{Factors affecting expressivity.}
\begin{itemize}
    \item Depth and width
    \item Choice of activation function
    \item Connectivity structure
\end{itemize}

\medskip

\noindent
\textbf{Piecewise linear networks.}  
ReLU networks represent functions composed of many linear regions.

\medskip

\noindent
\textbf{Geometric viewpoint.}  
Each layer partitions the input space, increasing the complexity of the representation.

\medskip

\noindent
\textbf{Tradeoffs.}
\begin{itemize}
    \item High expressivity $\rightarrow$ better approximation
    \item But increased risk of overfitting and optimization difficulty
\end{itemize}

\medskip

\noindent
\textbf{Implicit constraints.}  
Architectures restrict the class of functions that can be represented, providing inductive bias.

\medskip

\noindent
\textbf{Insight.}  
Expressivity alone is not sufficient; effective learning requires a balance between representation power and trainability.

\section{A Unifying Perspective}

Activation functions and architectures jointly determine:

\begin{itemize}
    \item how information flows through the network,
    \item how complex functions can be represented,
    \item how easily the model can be trained.
\end{itemize}

\medskip

\noindent
\textbf{Core components.}
\begin{itemize}
    \item Nonlinear activations
    \item Layer connectivity
    \item Architectural patterns
\end{itemize}

\medskip

\noindent
\textbf{Key insight.}  
The design of neural architectures is a central aspect of modern machine learning, combining mathematical structure with empirical performance.

\section{Outlook}

This chapter examined the role of activation functions and architectures:
\begin{itemize}
    \item nonlinear activations enabling expressive models,
    \item architectural patterns shaping learning,
    \item tradeoffs between depth and width,
    \item factors influencing expressivity.
\end{itemize}

\medskip

In the next chapter, we study normalization methods, which play a key role in stabilizing and accelerating training.

\medskip

\noindent
\textbf{Guiding principle.}  
The effectiveness of deep learning systems depends not only on their size, but on how they are structured.